% MyHatTileStuff/surrounding.c — Detailed Technical Explanation
% Compile with: pdflatex -shell-escape surrounding.tex
%
\documentclass[10pt,a4paper]{article}

%% ---------- Packages ----------
\usepackage[T1]{fontenc}
\usepackage[utf8]{inputenc}
\usepackage{geometry}
\geometry{margin=2cm, top=2.2cm, bottom=2.5cm}

\usepackage{microtype}
\usepackage{parskip}
\usepackage{xcolor}
\usepackage{booktabs}
\usepackage{amsmath,amssymb}
\usepackage{graphicx}
\usepackage{hyperref}
\usepackage{enumitem}
\setlist{noitemsep, topsep=2pt}
\usepackage{needspace}

%% minted — syntax highlighting (requires -shell-escape and pygments)
\usepackage{minted}
\setminted{
  fontsize=\footnotesize,
  baselinestretch=1.1,
  breaklines=true,
  linenos=false,
  frame=none,
  bgcolor=codebg
}
\definecolor{codebg}{HTML}{F5F5F5}

%% tcolorbox — side-by-side description/code panels
\usepackage[most]{tcolorbox}
\tcbuselibrary{minted,skins,breakable}

%% Colours
\definecolor{titlecolor}{HTML}{1A4D8F}
\definecolor{headercolor}{HTML}{2E6DB4}
\definecolor{rulecolor}{HTML}{AECBEF}
\definecolor{hatred}{HTML}{E74C3C}
\definecolor{hatblue}{HTML}{A8D8EA}

%% Side-by-side box style
\tcbset{
  dualbox/.style={
    enhanced,
    colframe=rulecolor,
    colback=white,
    fonttitle=\bfseries\small,
    breakable,
    boxrule=0.6pt,
    arc=3pt,
    sidebyside,
    sidebyside align=top seam,
    sidebyside gap=6mm,
    lefthand ratio=0.46,
    before skip=8pt,
    after skip=8pt,
  }
}

%% Section colours
\usepackage{titlesec}
\titleformat{\section}{\large\bfseries\color{titlecolor}}{}{0em}{}[\nobreak\vspace{-2pt}\color{rulecolor}\hrule\nobreak\vspace{4pt}]
\titleformat{\subsection}{\normalsize\bfseries\color{headercolor}}{}{0em}{}

%% Keep headings attached to the content that follows them
\newcommand{\subsectionbreak}{\nobreak}
\let\polysection\section
\renewcommand\section{\needspace{0.16\textheight}\polysection}
\let\polysubsection\subsection
\renewcommand\subsection{\needspace{0.12\textheight}\polysubsection}

%% Convenient macros
\newcommand{\code}[1]{\texttt{\small #1}}
\newcommand{\file}[1]{\texttt{\footnotesize #1}}

%% ---------- Document ----------
\begin{document}

\begin{center}
  {\Huge\bfseries\color{titlecolor} surrounding.c}\\[4pt]
  {\large\color{headercolor} Enumeration of Hat-Tile Surroundings of the Anti-Hat}\\[6pt]
  {\small \file{MyHatTileStuff/surrounding.c} — Technical Reference}\\[2pt]
  {\small\color{gray} \today}
\end{center}

\vspace{6pt}
\noindent\rule{\linewidth}{1pt}
\vspace{4pt}

% -----------------------------------------------------------------------
\section{Purpose}

\code{surrounding} enumerates all distinct ways to place unreflected hat
tiles around a fixed \emph{anti-hat} center tile so that:
\begin{itemize}
  \item every surrounding tile shares at least one edge directly with
    the center (direct contact, not just corner or proximity),
  \item the union of all tiles — center plus surrounding — has no holes
    (no enclosed empty regions), and
  \item at least \(n_{\min}\) surrounding tiles are present
    (default: 5; controllable via \code{--min}).
\end{itemize}

Two configurations are considered the \emph{same} if one can be obtained
from the other by a symmetry of the tetrille lattice (any combination of
rotations and reflections).  The result is a set of symmetry classes
visualised as an SVG grid.

The motivation is to characterise which local arrangements of regular
(unreflected) hat tiles can appear around an anti-hat in a valid hat
tiling, supporting the broader research goal of PolyformTown.

% -----------------------------------------------------------------------
\section{Background: Hat Tile and Anti-Hat}

\subsection{The Hat Tile}

The hat is a 13-edge polygon on the \textbf{tetrille} lattice, first
described in 2023 as the first known aperiodic monotile.  The tetrille
lattice is a kite-and-dart decomposition of the plane with three vertex
valences: $v \in \{3, 4, 6\}$.  Coordinates are integer triples
$(v, x, y)$ where $v$ carries the valence tag; cross-tag arithmetic is
impossible without explicit lifting formulas.

The hat tile has 14 vertices in its boundary cycle (the 13-edge polygon
closes back to the start), using all three valence systems.  Its
\file{hat.tile} DSL entry begins:
\begin{minted}{text}
name: hat
lattice: tetrille
constants:
r = sqrt(3)
cycle:
6  0  0   4 -1  0   3 -1  0
4 -1 -1   6  0 -1   4  0 -3
3 -1 -2   4  1 -3   3  0 -2
4  1 -2   6  1 -1   4  2 -1
3  1 -1   4  1 -1
\end{minted}

\subsection{Chirality: Hat vs.\ Anti-Hat}

The hat tile is \emph{chiral}: it and its mirror image (the
\textbf{anti-hat}) are geometrically distinct and cannot be superimposed
by any rotation.  Together they tile the plane aperiodically in a fixed
ratio of approximately $3{:}1$ (hats to anti-hats).

The tetrille lattice's symmetry group has 12 elements: transforms
$t = 0\ldots 5$ are pure rotations; $t = 6\ldots 11$ are a fixed
reflection followed by a rotation.  Applying $t = 6$ (pure reflection)
to the base hat cycle produces the anti-hat.

\code{tile\_build\_variants()} computes all 12 symmetry images of the
base cycle and deduplicates them.  For the hat tile this yields exactly
12 distinct variants (6 hat orientations and 6 anti-hat orientations)
because neither chirality has any non-trivial rotational symmetry.

% -----------------------------------------------------------------------
\section{Data Structures}

\subsection{Core Types}

\begin{tcolorbox}[dualbox, title={TileCyc and Config}]
  \textbf{TileCyc} stores one tile's boundary cycle compactly enough
  for SVG rendering without retaining a full \code{Cycle} structure.
  It records the vertex array, its length, and a chirality flag
  (retained for the centre tile even though surrounding tiles are always
  unreflected hats).

  \textbf{Config} records one complete configuration: the count of
  surrounding tiles and their \code{TileCyc} entries.  The centre is
  always \code{g\_center} (implicit), so it is not stored in the config.
\tcblower
\begin{minted}{c}
#define HAT_NMAX  20   /* 14 verts + margin */
#define MAX_SUR   18   /* DFS depth cap     */
#define MAX_CFGS  5000 /* output cap        */

typedef struct {
    int   n;
    Coord v[HAT_NMAX];
    int   reflected; /* 1=anti-hat, 0=hat */
} TileCyc;

typedef struct {
    int     sur_count;
    TileCyc sur[MAX_SUR];
} Config;
\end{minted}
\end{tcolorbox}

\subsection{Global State}

\begin{tcolorbox}[dualbox, title={Global Variables}]
  The program keeps all mutable state in file-scope statics so the
  DFS function signature stays clean.  The most critical globals are
  the per-depth aggregate array \code{g\_agg} and the DFS tile stack
  \code{g\_path}.

  \medskip
  \begin{description}[leftmargin=0pt, style=nextline]
    \item[\code{g\_tile}]
      The loaded tile (hat), including all 12 pre-computed variants.
    \item[\code{g\_refl[vi]}]
      Chirality flag: 1 if variant \code{vi} is an anti-hat.
    \item[\code{g\_center}]
      The anti-hat centre cycle (fixed, never changes).
    \item[\code{g\_cedges[\,]}, \code{g\_cedge\_n}]
      The 14 directed boundary edges of the centre tile; used to
      identify which aggregate boundary edges are eligible for
      attaching a surrounding tile.
    \item[\code{g\_seen}]
      A \code{HashTable} of canonical aggregate \code{Poly} shapes;
      each canonical state is visited at most once.
    \item[\code{g\_cfgs[\,]}, \code{g\_cfg\_n}]
      The collected output configurations.
  \end{description}
\tcblower
\begin{minted}{c}
static Tile      g_tile;
static int       g_refl[MAX_VARIANTS];
static Cycle     g_center;
static Edge      g_cedges[32];
static int       g_cedge_n;
static HashTable g_seen;
static Config    g_cfgs[MAX_CFGS];
static int       g_cfg_n;
static int       g_min_sur;
static int       g_maximal_only;

/* DFS tile stack: g_path[d] = tile placed at depth d */
static TileCyc   g_path[MAX_SUR];

/* Per-depth aggregate polys (see Section 5) */
static Poly      g_agg[MAX_SUR + 2];

/* Scratch temporaries (reused, never aliased) */
static Poly      g_canon_tmp;
static Cycle     g_aligned_tmp;
\end{minted}
\end{tcolorbox}

% -----------------------------------------------------------------------
\section{Initialisation}

\subsection{Loading and Variant Classification}

\begin{tcolorbox}[dualbox, title={Chirality Identification — \code{identify\_reflected()}}]
  After \code{tile\_load()} populates \code{g\_tile}, the program must
  determine which of the 12 variants are unreflected hats and which are
  anti-hats.

  The strategy: for each variant, apply all six pure-rotation transforms
  ($t = 0\ldots 5$) to the \emph{base} cycle and check if any result is
  coordinate-equal to that variant.  A variant reachable by a pure
  rotation is a hat; all others are anti-hats.

  Cycles are normalised (position and start vertex) before comparison so
  that coordinate equality is well-defined regardless of the lattice
  origin or cycle winding offset.
\tcblower
\begin{minted}{c}
static void identify_reflected(void) {
  for (int vi = 0; vi < g_tile.variant_count; vi++) {
    g_refl[vi] = 1; /* assume anti-hat */
    for (int t = 0; t < 6; t++) {
      Cycle tmp;
      cycle_transform_lattice(
          &g_tile.base, &tmp,
          g_tile.lattice, t);
      if (cycle_signed_area2(&tmp, g_tile.lattice) < 0)
          cycle_reverse(&tmp);
      cycle_normalize_position(&tmp, g_tile.lattice);
      /* coordinate-equal → it is a pure rotation */
      if (!cycle_less(&tmp, &g_tile.variants[vi]) &&
          !cycle_less(&g_tile.variants[vi], &tmp)) {
          g_refl[vi] = 0; /* it's a hat */
          break;
      }
    }
  }
}
/* Result for hat.tile: 6 hats (g_refl=0),
                        6 anti-hats (g_refl=1) */
\end{minted}
\end{tcolorbox}

\subsection{Centre Construction}

\begin{tcolorbox}[dualbox, title={Anti-Hat Centre and Edge Extraction}]
  The centre tile is constructed by applying transform $t = 6$ (the
  tetrille reflection) to the base hat cycle, normalising orientation
  (ensure CCW) and position.  This produces the canonical anti-hat
  in a fixed, reproducible coordinate frame.

  All 14 directed boundary edges of the centre are stored in
  \code{g\_cedges[]} at startup.  During DFS only aggregate boundary
  edges that appear in this set are eligible attachment sites,
  enforcing the direct-contact requirement.
\tcblower
\begin{minted}{c}
/* Build the anti-hat centre */
{
  Cycle tmp;
  cycle_transform_lattice(
      &g_tile.base, &tmp,
      g_tile.lattice, 6);      /* t=6 = pure reflection */
  if (cycle_signed_area2(&tmp, g_tile.lattice) < 0)
      cycle_reverse(&tmp);      /* ensure CCW */
  cycle_normalize_position(&tmp, g_tile.lattice);
  g_center = tmp;
}

/* Seed aggregate: centre alone */
g_agg[0].cycle_count = 1;
g_agg[0].cycles[0]   = g_center;

/* Record the 14 centre boundary edges */
{
  Edge tmp[64];
  g_cedge_n = build_boundary_edges(&g_agg[0], tmp);
  memcpy(g_cedges, tmp,
      (size_t)g_cedge_n * sizeof(Edge));
}
\end{minted}
\end{tcolorbox}

% -----------------------------------------------------------------------
\section{Depth-First Search}

\subsection{Overview}

The enumeration proceeds by depth-first search over sequences of tile
placements.  The DFS state at depth $d$ is the aggregate polygon
\code{g\_agg[d]}, which is the union of the centre tile plus the $d$
surrounding tiles placed so far.  At each depth the DFS:

\begin{enumerate}
  \item canonicalises \code{g\_agg[d]} and inserts it into
    \code{g\_seen}; if already seen, returns immediately (dedup),
  \item enumerates all valid one-tile extensions subject to the two
    constraints (see Section~6),
  \item recurses on each extension at depth $d + 1$,
  \item after all extensions are tried, records the current state
    as a finished configuration if it qualifies.
\end{enumerate}

\subsection{Per-Depth Aggregate Stack}

\begin{tcolorbox}[dualbox, title={Stack Layout — \code{g\_agg[0..MAX\_SUR+1]}}]
  A \code{Poly} struct is large
  ($\approx 144\,\text{KB}$ for \code{MAX\_CYCLES=32}, \code{MAX\_VERTS=384}).
  Placing one on the C call stack at each recursive level would risk a
  stack overflow at depth 18.

  The solution is a \emph{pre-allocated global array} indexed by depth.
  \code{g\_agg[d]} is the aggregate at depth $d$.  When the DFS at
  depth $d$ calls \code{try\_attach\_tile\_poly\_ex}, it passes
  \code{\&g\_agg[d+1]} as the output buffer.  The recursive call at
  depth $d{+}1$ then reads \code{g\_agg[d+1]} as its starting state
  and writes to \code{g\_agg[d+2]}, and so on.

  Because a level $d$ only modifies \code{g\_agg[d+1]}, \code{g\_agg[d]}
  remains valid across the recursive call.  Backtracking is free:
  the next extension at depth $d$ simply overwrites \code{g\_agg[d+1]}
  with a different placement.
\tcblower
\begin{minted}{text}
g_agg[0]  = centre tile only       (depth 0, set in main)
g_agg[1]  = centre + tile[0]       (written at depth 0)
g_agg[2]  = centre + tile[0,1]     (written at depth 1)
   ...
g_agg[d]  = centre + tile[0..d-1]
g_agg[d+1]= candidate new state    (written, then read
                                    at depth d+1)
\end{minted}
\vspace{6pt}

\begin{minted}{c}
/* DFS writes to g_agg[depth+1] before recursing */
if (!try_attach_tile_poly_ex(
      agg,                    /* g_agg[depth]   */
      &g_tile.variants[vi],
      g_tile.lattice,
      be, te,
      &g_agg[depth + 1],     /* output: next level */
      &g_aligned_tmp))
    continue;
dfs(depth + 1);
/* On return, g_agg[depth] is unchanged;
   g_agg[depth+1] will be overwritten by next
   iteration                                     */
\end{minted}
\end{tcolorbox}

\subsection{Deduplication}

\begin{tcolorbox}[dualbox, title={State Deduplication via Canonical Hash}]
  Many different tile-insertion \emph{sequences} produce the same
  aggregate \emph{shape}.  Without deduplication the DFS would explore
  exponentially many paths to the same state.

  At the top of \code{dfs(depth)}, \emph{before} any extension is
  tried, the aggregate \code{g\_agg[depth]} is canonicalised under all
  12 tetrille symmetries and the result is inserted into \code{g\_seen}.
  If \code{hash\_insert} returns 0 (already present), the call returns
  immediately.

  This is the standard ``canonical form at the root'' trick: since the
  canonical form is the same regardless of insertion order, each
  equivalence class of aggregate shapes is explored exactly once.
\tcblower
\begin{minted}{c}
static void dfs(int depth) {
  /* Canonicalise and dedup */
  poly_hash_key_lattice(
      &g_agg[depth],
      g_tile.lattice,
      &g_canon_tmp);          /* 12-transform minimum */
  if (!hash_insert(&g_seen, &g_canon_tmp))
      return;                 /* already explored */

  if (depth >= MAX_SUR) {
      if (depth >= g_min_sur) store_config(depth);
      return;
  }
  /* ... extension loop ... */
}

/* poly_hash_key_lattice applies all 12 tetrille
   transforms, normalises each image, and returns
   the lexicographically smallest one.  Two polys
   get the same canonical form iff they belong to
   the same symmetry class.                        */
\end{minted}
\end{tcolorbox}

\subsection{Extension Loop}

\begin{tcolorbox}[dualbox, title={Tile Attachment — Extension Loop}]
  For each boundary edge \code{be} of the current aggregate:
  \begin{enumerate}
    \item Skip \code{be} if it is not a centre edge
      (\code{is\_center\_edge()} — ensures direct contact).
    \item For each unreflected hat variant \code{vi}
      (anti-hats skipped via \code{g\_refl[vi]}):
    \item For each edge index \code{te} of variant \code{vi}:
    \item Call \code{try\_attach\_tile\_poly\_ex} to attempt
      placing variant \code{vi} with its edge \code{te}
      antiparallel to aggregate edge \code{be}.
    \item If the result has holes (\code{cycle\_count > 1}), skip.
    \item Otherwise record the placed tile in \code{g\_path[depth]}
      and recurse.
  \end{enumerate}
  The flag \code{extended} tracks whether any valid extension was found,
  used later for the maximal-mode filter.
\tcblower
\begin{minted}{c}
const Poly *agg = &g_agg[depth];
Edge edges[256];
int ec = build_boundary_edges(agg, edges);

int extended = 0;
for (int be = 0; be < ec; be++) {
  if (!is_center_edge(edges[be])) continue;

  for (int vi = 0; vi < g_tile.variant_count; vi++) {
    if (g_refl[vi]) continue;  /* hats only */

    for (int te = 0; te < g_tile.variants[vi].n; te++) {
      if (!try_attach_tile_poly_ex(
              agg, &g_tile.variants[vi],
              g_tile.lattice, be, te,
              &g_agg[depth + 1], &g_aligned_tmp))
          continue;

      /* Constraint: no holes */
      if (g_agg[depth + 1].cycle_count > 1)
          continue;

      extended = 1;
      /* Copy placed tile into path stack */
      int n = g_aligned_tmp.n;
      if (n > HAT_NMAX) n = HAT_NMAX;
      g_path[depth].n         = n;
      g_path[depth].reflected = g_refl[vi]; /* =0 */
      for (int i = 0; i < n; i++)
          g_path[depth].v[i] = g_aligned_tmp.v[i];

      dfs(depth + 1);
    }
  }
}
\end{minted}
\end{tcolorbox}

\subsection{Configuration Recording}

After the extension loop, the current depth-$d$ state is recorded if it
satisfies the output criteria:

\begin{equation*}
  \texttt{depth} \ge g\_min\_sur
  \quad\text{and}\quad
  (\lnot\texttt{g\_maximal\_only} \;\lor\; \lnot\texttt{extended}).
\end{equation*}

In default mode every state with at least \(g\_min\_sur\) tiles is
stored.  In \code{--maximal} mode only states from which no further
valid hat tile can be placed (touching the centre, without creating a
hole) are stored.

% -----------------------------------------------------------------------
\section{Key Constraints}

\subsection{Direct Contact with the Centre}

Every surrounding tile must share a boundary edge with the centre
anti-hat, not merely be adjacent to a previously placed surrounding
tile.

This is enforced by \code{is\_center\_edge(e)}: it returns 1 only if
the directed edge $e$ appears in the pre-recorded centre edge set
\code{g\_cedges[]}.  Because tiles are attached to \emph{boundary}
edges of the growing aggregate, and because a centre edge disappears
from the aggregate boundary once a tile is attached there, this
filter ensures that every tile in the DFS is placed directly against
the centre.

\begin{minted}{c}
static int is_center_edge(Edge e) {
    for (int i = 0; i < g_cedge_n; i++)
        if (coord_eq(e.a, g_cedges[i].a) &&
            coord_eq(e.b, g_cedges[i].b))
            return 1;
    return 0;
}
\end{minted}

\noindent Note: \code{coord\_eq} checks all three fields $(v, x, y)$,
so the comparison is exact and tag-safe.

\subsection{Unreflected Hats Only}

Anti-hat tiles are excluded from the surrounding by the single guard:
\begin{minted}{c}
if (g_refl[vi]) continue;
\end{minted}
inside the variant loop.  Only the six unreflected hat variants
($t = 0\ldots 5$) are ever tried as surrounding tiles.

\subsection{No-Hole Filter}

The \code{Poly} struct produced by \code{try\_attach\_tile\_poly\_ex}
can have multiple boundary cycles: \code{cycles[0]} is always the outer
boundary (CCW), and \code{cycles[1\ldots]} are hole boundaries (CW).
A \code{cycle\_count} greater than 1 signals that the placement would
enclose an empty region.

\begin{minted}{c}
if (g_agg[depth + 1].cycle_count > 1) continue;
\end{minted}

This check is applied \emph{immediately after} attachment, before
recursing.  Holes are monotone: once created, no additional tiles on the
exterior can eliminate them.  Early pruning therefore cuts entire
subtrees, not just individual leaf configurations.

% -----------------------------------------------------------------------
\section{Tile Attachment Mechanics}

\code{try\_attach\_tile\_poly\_ex} is the geometric core from
\file{src/core/attach.c}.  Given aggregate \code{agg}, a tile variant
cycle \code{tv}, boundary edge index \code{be}, and tile edge index
\code{te}, it:

\begin{enumerate}
  \item \textbf{Aligns} the tile so that the directed tile edge
    \code{tv.v[te]} $\to$ \code{tv.v[te+1]}
    is antiparallel and coincident with the directed aggregate edge
    \code{edges[be]}.  The required translation is computed in the
    6-system (the coarsest integer lattice that contains all three
    valence systems) to avoid cross-tag arithmetic.
  \item \textbf{Collects} all directed edges from both the aggregate
    and the aligned tile into a single flat list.
  \item \textbf{Cancels} pairs of antiparallel coincident edges (the
    shared edges that disappear when two tiles are joined).  Parallel
    coincident edges indicate overlap and abort with \code{ATTACH\_STATUS\_GEOMETRY}.
  \item \textbf{Walks} the remaining directed edges to extract closed
    boundary cycles, classifying each by signed area:
    $A_2 > 0$ (CCW) $\Rightarrow$ outer; $A_2 < 0$ (CW) $\Rightarrow$ hole.
  \item Writes the aggregate result to \code{g\_agg[depth+1]} and the
    aligned tile's own cycle to \code{g\_aligned\_tmp}.
\end{enumerate}

The \code{g\_aligned\_tmp} output is the tile in the shared coordinate
frame — its vertex array is copied into \code{g\_path[depth]} for
later SVG rendering.

% -----------------------------------------------------------------------
\section{Canonicalisation}

\code{poly\_hash\_key\_lattice(src, lattice, key)} reduces a \code{Poly}
to its canonical representative under the 12-element tetrille symmetry
group ($D_6$ acting on the tagged integer lattice):

\begin{enumerate}
  \item For each transform $t \in \{0, \ldots, 11\}$:
    \begin{enumerate}
      \item Apply \code{poly\_transform\_lattice(src, \&tmp, lattice, t)}.
      \item Translate \code{tmp} so its lexicographically smallest
        vertex is at the system-appropriate origin
        (\code{poly\_normalize\_position}).
      \item Rotate each cycle to start at its own lexicographically
        smallest vertex (\code{cycle\_canonicalize\_shift}).
      \item Sort hole cycles lexicographically.
    \end{enumerate}
  \item Select the transform image \code{tmp} that is lexicographically
    smallest overall as the canonical key.
\end{enumerate}

Two \code{Poly} shapes hash to the same key if and only if they are
related by a tetrille symmetry, which is exactly the condition for them
being \emph{equivalent} surroundings.  The hash table \code{g\_seen}
uses this key; \code{hash\_insert} returns 0 on collision, preventing
re-exploration of equivalent aggregate states.

\medskip
\noindent\textbf{Consequence for DFS.}  The dedup check fires at the
\emph{start} of each DFS call, before any new tile is placed.  This
means the state \emph{before extension} is tested — equivalently, the
set of all partial placements up to and including depth $d$ is
represented by a single canonical aggregate rather than by the
(exponential) set of tile-insertion sequences that produce it.

% -----------------------------------------------------------------------
\section{SVG Rendering}

\subsection{Coordinate Conversion}

Tetrille lattice coordinates $(v, x, y)$ are mapped to real screen
coordinates using the basis matrices from \file{hat.tile}.  The
$\sqrt{3}$ factor is pre-computed once as \code{g\_r3 = sqrt(3.0)}.

\[
  \begin{pmatrix} x_\text{screen} \\ y_\text{screen} \end{pmatrix}
  =
  \begin{cases}
    \begin{pmatrix} \tfrac{\sqrt{3}}{2}(x+y) \\ \tfrac{1}{2}(y-x) \end{pmatrix}
    & v = 6 \\[6pt]
    \begin{pmatrix} \tfrac{\sqrt{3}}{4}(x+y) \\ \tfrac{1}{4}(y-x) \end{pmatrix}
    & v = 4 \\[6pt]
    \begin{pmatrix} \tfrac{1}{\sqrt{3}}x + \tfrac{1}{2\sqrt{3}}y \\ \tfrac{1}{2}y \end{pmatrix}
    & v = 3
  \end{cases}
\]

SVG $y$ increases downward while mathematical $y$ increases upward, so
the screen $y$-coordinate is negated: a vertex at logical $(x_s, y_s)$
maps to SVG position $(o_x + \texttt{SVG\_UNIT}\cdot x_s,\;
o_y - \texttt{SVG\_UNIT}\cdot y_s)$ where $(o_x, o_y)$ is the cell
origin offset.

\subsection{Layout}

Configurations are sorted by surrounding tile count and arranged in a
grid whose aspect ratio approximates $16{:}9$:
\begin{itemize}
  \item A single global bounding box is computed over all tile vertices
    in all configurations, ensuring all cells have the same scale.
  \item Each cell receives a label ``\code{n=\textit{k}}'' below the
    drawing.
  \item Surrounding tiles are drawn first (behind the centre); the
    centre anti-hat is drawn on top.
\end{itemize}

Colour coding:
\begin{itemize}
  \item \textbf{\color{hatred}Coral} (\texttt{\#e74c3c}): centre
    anti-hat.
  \item \textbf{\color{hatblue}Light blue} (\texttt{\#a8d8ea}):
    surrounding (unreflected) hat tiles.
\end{itemize}

% -----------------------------------------------------------------------
\section{Results}

Running \code{./bin/surrounding} with the default hat tile produces:

\medskip
\begin{tabular}{@{}lll@{}}
\toprule
\textbf{Mode} & \textbf{Configurations} & \textbf{Distribution} \\
\midrule
Default (\code{--min 4})  & 108 & all $n = 4$ \\
\code{--maximal}          & 108 & all $n = 4$ \\
\bottomrule
\end{tabular}

\medskip
\noindent The identical counts in both modes reveal that every valid
4-hat surrounding is already maximal: no 5th unreflected hat tile can
be attached to any remaining centre edge without either overlapping an
existing tile or creating a hole.  This is a non-trivial geometric
constraint arising from the specific shape of the anti-hat and the
chirality restriction.

Diagnostic output printed to \code{stderr}:
\begin{minted}{text}
=== Hat Tile Surroundings ===
Tile file:       tiles/hat.tile
Variants:        12 total (6 hat, 6 anti-hat)
Center edges:    14
Min surrounding: 4
Mode:            maximal only

Configurations:  108
States visited:  337
\end{minted}

% -----------------------------------------------------------------------
\section{Usage}

\begin{minted}{bash}
# Build
make surrounding

# All configurations with >= 4 surrounding hats
./bin/surrounding

# Only maximal configurations
./bin/surrounding --maximal

# Custom minimum, custom output file
./bin/surrounding --min 3 --out my.svg

# Different tile file
./bin/surrounding --tile tiles/hat.tile --out surrounding.svg
\end{minted}

The output SVG is written to \code{surrounding.svg} (or the path given
by \code{--out}).  It is self-contained and can be opened in any
standards-compliant SVG viewer or web browser.

% -----------------------------------------------------------------------
\section{Architecture Summary}

\begin{center}
\begin{tcolorbox}[
  colframe=rulecolor, colback=white, boxrule=0.6pt, arc=3pt,
  width=0.96\linewidth
]
\begin{minted}{text}
tile_load("tiles/hat.tile")
  +-- tile_build_variants() -> 6 hat + 6 anti-hat variants
         |
         +-- identify_reflected() -> g_refl[vi]
         |
         +-- cycle_transform_lattice(base, t=6) -> g_center (anti-hat)
         |
         +-- build_boundary_edges(g_agg[0]) -> g_cedges[14]

dfs(depth)
  +-- poly_hash_key_lattice(g_agg[depth]) -> g_canon_tmp
  |    +-- if seen: return (dedup)
  |
  +-- build_boundary_edges(g_agg[depth]) -> edges[]
  |
  +-- for each edge be:
       +-- if is_center_edge(be):       [direct contact]
            +-- for each variant vi (g_refl[vi]==0): [hats only]
                 +-- for each tile edge te:
                      +-- try_attach_tile_poly_ex(...)
                           +-- if fail: skip
                           +-- if cycle_count > 1: skip [no holes]
                           +-- copy g_aligned_tmp -> g_path[depth]
                                dfs(depth + 1)

  +-- if depth >= g_min_sur && (all||!extended): store_config(depth)

emit_svg() -> surrounding.svg
\end{minted}
\end{tcolorbox}
\end{center}

\vspace{4pt}
\noindent\rule{\linewidth}{0.6pt}
\vspace{2pt}

{\small\color{gray}
\file{MyHatTileStuff/surrounding.c} — part of PolyformTown.
Compile this document with \texttt{pdflatex -shell-escape surrounding.tex}
(requires \texttt{pygments} for \texttt{minted}).
}

\end{document}
